\chapter{Data reconstruction}\label{app:data}
\section{Sources and alignment}
We archive the five course inputs and align their date keys on a continuous
UTC calendar from 1 January 2020 to 17 September 2026. The original fuel
total is checked against the fuel breakdown. When the fuel data are
reconstructed, their sum replaces the reported total. The wind-farm file is
a simulation. We retain its ten zero-output days and exclude its 18 September
2026 row from the common calendar. Within that calendar the simulation
values are preserved. The original files omit 30 generation dates,
one demand date and 30 price dates. The repairs use Elexon's
\href{https://bmrs.elexon.co.uk/api-documentation/endpoint/generation/actual/per-type}{AGPT},
\href{https://bmrs.elexon.co.uk/api-documentation/endpoint/datasets/FUELHH}{FUELHH},
\href{https://bmrs.elexon.co.uk/api-documentation/endpoint/demand/outturn/summary}{VD}
and \href{https://bmrs.elexon.co.uk/api-documentation/endpoint/datasets/MID}{MID}
archives, and \href{https://www.neso.energy/data-portal/historic-demand-data}
{NESO historical demand and embedded generation} \cite{repair}.

\section{Generation}
\textbf{Missing and zero values.}
The simultaneous zeros in eight fuel categories on 107 course dates are
inconsistent with contemporaneous Elexon generation records. We use
complete Elexon Actual Generation per Production Type (AGPT) data where
available and reconstruct the eight categories from Elexon half-hourly
fuel outturn (FUELHH) where needed. Gas combines combined- and
open-cycle generation. Pumped storage counts positive output only.
Other coincident partial zeros are flagged separately from normal
fuel-specific zeros. For the 42 zero-solar dates, we prefer complete AGPT
records, interpolate short gaps and estimate heavily missing days from
calibrated National Energy System Operator (NESO) embedded solar data.
One date belongs to both zero-value groups. Validation checks that all
107 fuel dates and 42 solar dates were repaired.

\textbf{Incomplete coverage and export differences.}
Seven course days reproduce the midnight query endpoint. We reaggregate
the complete UTC days. The current AGPT archive contains 264 incomplete
days, and
136 days with simultaneous zero periods. For short gaps of at most six
half-hours, we interpolate from valid observations on the same day, using
the nearest observation at a day boundary. For larger gaps, the eight
non-wind/solar categories use FUELHH and renewable components use
calibrated alternative data. Solar is scaled to AGPT by year. Wind
uses a local fit based on FUELHH wind and NESO embedded wind, with the
nearest 60 complete training days in the same year and every fifth day
held out. We rebuild the 87 differences between the 2026 extract and
the current source, including 61 dates outside the initial 107-date repair.

\newpage
\section{Demand and market index prices}
\textbf{Demand.}
The original demand data reproduce Elexon demand outturn (VD) query
endpoints on 349 Tuesdays and the final file date. Requerying the full
day removes this sampling error. The retrieved VD output has a
five-minute grid, with 116 days containing fewer than 288 records and
42 days having readings below 10 GW. We keep valid readings and fill
gaps with contemporaneous NESO transmission system demand plus a local
median VD-to-NESO offset. Nearby complete days supply the offset on two
dates with only one VD record. The method reconstructs 153 days. On
2 August 2022 the upstream feed has 17 consecutive readings at 636 MW.
Replacing those intervals and averaging the full day yields an estimated
28,744 MW.

\textbf{Prices.}
Elexon Market Index Data (MID) reproduce the 346 query-end prices and
the remaining 2,076 published daily values under the original rule.
We remove 115,571 provider records with zero reported traded volume,
then volume-weight active providers within each half-hour and average
those half-hours over the UTC day. Genuine zero or negative prices
with traded volume are retained. On 1 January 2020 the daily price
changes from \gbp 18.59 to \gbp 37.19/MWh. Seven nominally negative
course dates become positive on complete-day averaging. The repaired
series has two negative daily means. Short gaps are interpolated.
For 35 missing or severely incomplete price days, estimates use four
nearby observed days with the same weekday. We maintain both price
columns for compatibility and use the \gbp/MWh series in the model.

\section{Repair outcomes and verification}
The release contains 2,452 observations per series, a cell-level record
of original and replacement values, and flags for the method used.
Across the calendar, 436 dates use an estimate or interpolation,
323 use cross-source reconstruction, and 1,929 meet the strict
quality flag. The categories can overlap. Independent checks reconcile
fuel components to total generation, verify the price-column factor of
24 and recompute 2,372 prices with complete active-provider coverage.

Held-out comparisons measure the uncertainty of reconstruction. Annual
solar mean absolute error ranges from 42 to 182 MW, onshore wind from
363 to 592 MW and offshore wind from 453 to 832 MW. Hiding the first
17 demand intervals on 74 control days gives a daily mean absolute
error of 8.63 MW. Hiding a whole price day on 77 control days gives
\gbp 19.87/MWh mean absolute error. These validation errors describe
the held-out days. The archived inputs, evidence, audit and validation
records are identified in \cite{repair}.

\chapter{Market and equipment research}\label{app:research}
\section{Electricity trading and hydrogen offtake}
\textbf{Electricity.}
The course demand series represents a constructed market. Published
imbalance settlement applies to differences between contracted and metered
positions. A company-specific charge would require an actual delivery
commitment, settlement terms and an applicable price. We model sale at
the daily market price and report demand coverage separately. Policy B
is a dispatch scenario for testing priority.

\textbf{Hydrogen price.}
The sources examined provide distinct measures.
\href{https://www.eex.com/en/markets/hydrogen}{EEX HYDRIX} is a weekly
German price indication including delivery.
\href{https://www.icis.com/explore/resources/news/2022/09/27/10810133/icis-launches-europe-s-first-market-linked-renewable-hydrogen-assessments/}{ICIS}
publishes assessed values linked to production costs and power arrangements.
\href{https://www.argusmedia.com/en/solutions/products/argus-hydrogen-and-future-fuels}{Argus}
provides modelled costs, including UK proxies. The UK Hydrogen Allocation Round 1
(HAR1) reports a supported strike price, while the Low Carbon Hydrogen
Agreement distinguishes the buyer's achieved sales price from support
payments. We use the HAR1 reference as a constant netback scenario.
At the benchmark of \gbp 175/MWh in 2012 prices and 39.4 kWh/kg, the
converted value is \(175(39.4/1000)(108.7/96.1)=7.799\) in
2020 \gbp/kg. The scenario applies no additional non-electricity
deduction and assumes the buyer accepts each model's gas quality
\cite{har,cpi,basis,v6}.

\textbf{Hydrogen volume.}
The European Hydrogen Observatory's 2025 update reports about 569,024,
474,414 and 485,626 tonnes of UK consumption in 2022, 2023 and 2024.
We divide by the respective calendar days to obtain approximately
1,559, 1,300 and 1,327 t/day. The largest is a broad market case that
smooths aggregate consumption across the year and includes existing
captive use. The 1.56 t/day case is 0.1\% of that benchmark and
10 t/day is an illustrative buyer case. Actual offtake depends on
location, purity, delivery capacity and contract terms \cite{eho,v6}.

\section{Battery selection and cost basis}
Sungrow ST255CS-2H, Huawei LUNA215 and LUNA241 and Tesla Megapack
2XL are compatible with the two-hour reference. We use model-specific
power, nominal capacity and AC round-trip efficiency, splitting the latter
symmetrically between charge and discharge. The 20--80\% operating
band is a common modelling convention. It provides 60\% of nominal
capacity as usable stock. In the main case, stated AC efficiency
covers round-trip energy losses and we apply zero explicit daily
self-discharge \cite{brief,dea,v6}.

For the three smaller units, the Danish Energy Agency installed-cost
reference is scaled by energy capacity and converted to 2020 pounds.
The fixed annual budget is 2.9\% of capital. We also test the catalogue's
literal capacity-based operating-cost line as an alternative. For
Megapack, the \href{https://web.archive.org/web/20240106085048id_/https://www.tesla.com/megapack/design}{2024 US public estimate}
is US\$9.76m installed for five
2XL units. The model transfers the average price of US\$1.95m per
unit to single-unit decisions at 2020 exchange and price equivalence.
The 2026 equipment-only listing omits installation and tax. For a
lower-cost case we add a Lazard-based installation proxy. The Lazard
cost boundaries distinguish standard maintenance, warranty and capacity
augmentation. We adopt the US National Renewable Energy Laboratory
\href{https://atb.nrel.gov/electricity/2024/utility-scale_battery_storage}{2.5\%-of-capital fixed annual budget}
for the maintained-capacity case.
It includes the upkeep provision, so no separate augmentation bill is
added. The historical US estimate provides a scenario price for the
UK project. The available sources do not establish a bulk-discount
curve \cite{dea,basis,v6}.

\newpage
\section{Electrolyser selection and cost basis}
The eight models in Table~\ref{tab:equipment} use published manufacturer
ratings. We obtain actual input MW from the rated hourly output and
specific electricity consumption and use the product's minimum sustained
load. Equipment is selected in whole units. The wet-gas POSEIDON
reference would require a matching treatment and delivery cost for
comparison with the selected product grades \cite{brief,v6}.

For the main comparable case, the midpoint of the International Energy
Agency's 2024 non-China installation range of US\$2,000--2,600/kW
is converted to \gbp 1.46m per MW in 2020 prices. The UK production
cost study supplies fixed maintenance and absolute stack replacement
costs, normalised to AC input power. The alternative uses the UK cost
study's capital and fixed-operation package with a variable replacement
reserve per kg, in place of the scheduled replacement charges
\cite{iea,beis,basis,v6}.

\section{Financial conversion and operating horizon}
Monthly nominal electricity prices are multiplied by the ratio of the
2020 mean Consumer Prices Index to the index in the observation month.
The course battery reference of EUR 288,000/MWh corresponds to
\gbp 256,235/MWh at the 2020 exchange rate. The comparable US equipment
conversion factor is 0.635 in 2020 pounds per 2024 dollar. Annual
equivalent costs use a 5\% real discount rate, 15 years for batteries
and 30 years for electrolysers, with zero residual value. The
commissioning charge purchases the starting 20\% battery fill at
\gbp 78.13/MWh, allowing for charge efficiency. Subsequent stock is
measured above that level and returns to it at the end of the study.
Electricity infrastructure costs that are identical across investments
cancel against each policy's no-investment comparator \cite{cpi,basis,v6}.
